\begin{table}[h!]
\centering
\begin{tabular}{| p{0.22\textwidth} | p{0.38\textwidth} | p{0.38\textwidth} |}
\hline
\textbf{Jenis Risiko} 
& \textbf{Deskripsi Risiko} 
& \textbf{Mitigasi} \\
\hline

Risiko Data
& \textit{Dataset} tidak seimbang, terdapat \textit{noise} atau duplikasi.
& Menambah variasi data, melakukan pembersihan ketat, dan menerapkan teknik \textit{balancing} sederhana. \\
\hline

Risiko Model \& Komputasi
& Dimensi fitur besar menyebabkan pelatihan lama atau performa model rendah.
& Membatasi vocabulary TF-IDF, menyederhanakan kategori sumber, serta melakukan \textit{tuning} ringan. \\
\hline

Risiko Jadwal
& Keterlambatan implementasi atau perubahan arahan pembimbing.
& Menyusun \textit{timeline} rinci, melakukan konsultasi rutin, dan menjaga ruang lingkup tetap fokus. \\
\hline

Risiko Ketidaksesuaian Rencana
& Dataset kurang atau hasil model di bawah ekspektasi.
& Menyempitkan ruang lingkup Tugas Akhir atau menjadikan analisis keterbatasan sebagai kontribusi ilmiah. \\
\hline

\end{tabular}
\caption{Identifikasi risiko dan strategi mitigasi}
\label{tbl:risiko-mitigasi}
\end{table}
